\section{Related Work}
\label{sec:related}
This section relates AgentWorkShop to plant-facing industrial agents and to
production and automation abstractions for context, tool access, and
operational constraints. Table~\ref{tab:nearest} compares integration and
execution features described in selected industrial-agent architectures,
including systems with adjacent but distinct process-control objectives. The
comparison is focused rather than exhaustive; a feature marked as not
specified was not described in the cited study and is not inferred to be
absent.
Studies were included in this comparison when they describe an industrial
agent's production context, plant-facing operations, or an action, approval,
or evidence pathway. This criterion defines the table's scope; it is not a
comprehensive survey or a performance ranking.
Feature entries were coded from the cited study's full text only when the
feature was explicitly described; otherwise the cell is marked ``Not
specified.'' Evaluation descriptions are transcribed at the level reported
by each source, without normalizing sample units, success criteria, or
metrics across studies.

\subsection{Industrial Agents and Plant Integration}
Industrial-agent systems, including large language model (LLM)-based systems,
span end-to-end automation control, supervisory
services, programmable logic controller (PLC) interfaces, manufacturing
coordination, and operator-facing machine access. Xia et al. present an
LLM agent framework for end-to-end industrial automation control, using
structured prompts and event-driven context for production planning and
operation \cite{xia2025control}. Huang et al. present a multimodal
architecture whose stated focus is automated control generation for
industrial cyber-physical systems \cite{huang2026made}. Biskupovic et al.
study agent-based supervision for service-oriented industrial systems
\cite{biskupovic2025supervision};
Agents4PLC and InstructMPC address PLC code generation and human--LLM
context-aware control, respectively
\cite{liu2026agents4plc,wu2025instructmpc}. Manufacturing studies consider
dynamic resource-capability exploration \cite{lim2025llmenabled}, multi-agent
shop-floor coordination \cite{zhao2026shopfloor}, and digital-twin-supported
fault handling \cite{gill2025llmdt}. Natural-language machine access
\cite{hofmann2025nlc} and real-time LLM control of a robot-operated serial
line \cite{waseem2026pretrained} extend this work to operator interfaces and
robotic-line control. Together, these studies address different operational
contexts, interfaces, and allocations of responsibility between agents,
operators, and controllers.

The closest architecture-level comparison is Xia et al., who connect cognitive,
digital-twin, and physical-automation layers and organize twin information as
state snapshots, event messages, and plan sequences; their case studies
address robotic automation and process simulation
\cite{xia2025architecture}. Other nearby designs expose factory-module
operations through MCP with state-based tool filtering, plan constrained
manufacturing steps, or compose semantic recipes from OPC UA skills
\cite{koehle2026llmdriven,silva2025mcpmanufacturing,koehle2025semanticrecipes}.
These studies describe important parts of plant-facing integration. The
recent systems in Table~\ref{tab:nearest} also combine agent execution with
approval, industrial data, or audit records. Graul et al. ground proposals in
the Asset Administration Shell (AAS) and route approved actions through OPC UA
Methods, with signed approvals and audit logging
\cite{graul2026aas}. Schall separates advisory language-model agents from
process and instrumentation diagram (P\&ID)-grounded verification and
execution, and evaluates structural checks across five simulated process
scenarios \cite{schall2026safe}; its context includes supervisory control and
data acquisition (SCADA) and standard operating procedure (SOP) information.
Lin and Chen build a governed sensor--MES--ERP digital thread for closed-loop
monitoring and control, evaluated on real manufacturing
records \cite{lin2026agentic}. These works address asset-specific approval,
process validation, or closed-loop analytics. AgentWorkShop examines the
complementary software-integration question of whether distinct registered
task adapters and a direct-call surface can reach one named industrial
operation through a common handler while retaining route-appropriate records.
The evaluated contract covers production context, semantic I/O, configured
drivers, command receipts, and later process observations. Its route-specific
approval check also rechecks the line--product--recipe--run tuple immediately
before software dispatch. This is a scoped implementation contract, not a
claim that other systems lack equivalent mechanisms. The study tests the
contract with two registered adapters, direct REST, and two simulator
mappings; it does not establish general adapter interchangeability, reduced
integration effort, or superiority over another platform.
Table~\ref{tab:nearest} summarizes only
features explicitly described in the cited studies.
\begin{table*}[!t]
\centering
\caption{Feature-level comparison of selected plant-facing integration
studies. Entries summarize reported design scope and evaluation, not performance.}
\label{tab:nearest}
\small
\setlength{\tabcolsep}{3pt}
\begin{tabularx}{\textwidth}{@{}>{\raggedright\arraybackslash}p{1.45cm}>{\raggedright\arraybackslash}p{3.85cm}>{\raggedright\arraybackslash}p{2.7cm}>{\raggedright\arraybackslash}X@{}}
\toprule
System & Context and operation checks & Device-facing route & Reported evaluation \\
\midrule
Xia et al. \cite{xia2025architecture} & Cognitive, twin, and physical layers; snapshots, events, plans & Digital-twin/physical automation & Robotic automation and process-simulation cases \\
K{\"o}hle et al. \cite{koehle2026llmdriven} & Factory agents with injected live state; state-based tool filtering & MCP exposes OPC UA methods; MQTT coordination & Nine simulated production challenges \\
Vieira da Silva et al. \cite{silva2025mcpmanufacturing} & Constrained LLM planning; operation checks not specified & MCP-exposed manufacturing functions & Laboratory multi-step process \\
K{\"o}hle et al. \cite{koehle2025semanticrecipes} & AI-assisted recipes from OPC UA skills; checks not specified & OPC UA skills and virtual commissioning & Unity assembly-workstation proof of concept \\
Graul et al. \cite{graul2026aas} & AAS-grounded context; signed human approval & PLC-validated OPC UA Methods; immutable audit log & Two laboratory PoCs; single scripted runs \\
Schall \cite{schall2026safe} & P\&ID, SCADA, time-series, SOP; topology and rule checks & Advisory proposal with gated execution & Five simulated scenarios; 43 proposals and 50 LLM runs \\
Lin and Chen \cite{lin2026agentic} & Sensor--MES--ERP thread; SOP-, risk-, capability-aware gating & Edge monitoring and closed-loop actions & 25,275 real records; operational metrics \\
\midrule
This work & OMP and Pi registered task adapters plus direct calls; approval-bound context rechecked before recipe dispatch & Shared dispatcher; semantic nodes resolve to configured DAQ/DCW drivers & Three tested entry paths across two simulator mappings; deterministic Pi fixture; no attainment claim \\
\bottomrule
\end{tabularx}
\end{table*}

The operation checks complement safety-aware decision support
\cite{jiang2026agentic} but do not replace independent process protection.

\subsection{Production Context, Agent Interfaces, and Automation}
Industrial integration also depends on abstractions that are not specific to
agent architectures. The Model Context Protocol (MCP) specifies interoperable
tool discovery and call mechanisms \cite{mcp2024}; the operation-level scope,
policy, grant, and binding checks discussed here are AgentWorkShop's own
dispatcher rules, distinct from MCP transport authorization. In parallel,
ISA-88 and ISA-95 describe recipes and production operations
\cite{isa88_2010,isa95_2025}. Run-to-run methods connect process observations
to successive production runs \cite{moyne2000r2r}. At the agent level,
AgentSpec and CaMeL address tool behavior and capability-based information
flow \cite{wang2025agentspec,debenedetti2025camel}. Together, these works
provide complementary abstractions for production context, tool interaction,
and information flow. AgentWorkShop composes such abstractions with its task
runtime, operation checks, semantic I/O, and route-specific records; this
paper evaluates selected paths through that composition rather than proposing
a replacement for the underlying standards.
